\section{Related Work}
\label{sec:related}

Our work sits at the intersection of three recent threads: scepticism
about chain-of-thought on in-context tasks, the limits of LLM spatial
reasoning, and surface-repetition effects in prompts.

\para{When does chain-of-thought help?}
\citet{zheng2025curse} test 16 LLMs across nine pattern-induction
benchmarks and report a \emph{Curse of CoT}: direct prompting beats CoT
by 20\% on average and 42\% on symbolic-grid tasks, with the gap
widening as more demonstrations are added. The proposed mechanism is
that CoT injects noise via flawed rationales while increasing the
contextual distance between demonstrations and the answer, disrupting
the implicit ICL pathway. \citet{sprague2025tocot} provide a
complementary meta-analysis: across 20 datasets and 14 LLMs, CoT
substantially helps only on math and symbolic-execution tasks, with
95\% of CoT's gain on MMLU concentrated in questions containing the
``='' symbol. Unlike both lines, which treat pattern induction
holistically and conclude that CoT hurts, we hold the
geometry fixed and ask about \emph{addressing} a known grid rather than
inducing a rule; the per-query-type analysis we report (\secref{sec:results})
shows that the two pictures are consistent once one separates
scan-and-aggregate queries from random-access queries.

\para{LLM spatial reasoning and shape understanding.}
\citet{rudman2025polygons} show that multimodal LLMs fall to
near-chance accuracy on novel polygons and that perception, not
language, is the bottleneck: their VC-CoT intervention---visually
labelling vertices with letters and asking the model to enumerate
them---lifts GPT-4o from 7\% to 93\% on irregular-polygon side counts.
\citet{wang2025vdlm} introduce Primal Visual Description, a vertex/colour
JSON schema, and show that a text-only LLM consuming PVD beats GPT-4V on
shape comparison and maze tasks. \citet{wu2024vot} propose
Visualization-of-Thought (VoT) for text-grid navigation, re-rendering
the grid at each step. \citet{bai2025matrix} stress-test Claude 3.7
Sonnet (16k thinking tokens vs.\ no-thinking), GPT-4o, and GPT-4.1 on
five ASCII-grid spatial tasks and find accuracy drops up to 84\% as grid
size grows. We use the same family of text-grid inputs, but our axis
is the \emph{type of question} asked about a fixed grid rather than
\emph{grid size} or \emph{representation}; we show that under \COT the
question-type axis collapses to ceiling for both tested models at
$7{\times}7$.

\para{Repetition and in-context learning.}
\citet{yan2024repetitions} show that ten repetitions of any token
sequence push its predicted probability to ${\approx}1.0$, evidence of a
strong self-reinforcement effect that is purely surface-statistical.
\citet{liu2025sparks} show that LLMs pick up 2D-geometric structure
from a few in-context $(\text{point}, \text{value})$ pairs, and
\citet{zheng2024memorization} provide an axiomatic decomposition of
ICL into memorisation and reasoning components. Curse-of-CoT treats
demonstration count as a confound rather than a hypothesis variable.
We deliberately make $K$ an axis (\cf\ \secref{sec:method}), so that
the joint effect of reasoning and repetition can be read off the
$2{\times}2{\times}5$ design directly. To our knowledge no prior work
crosses verbal addressing of novel 2D shapes with both axes
simultaneously.

\para{Novelty in evaluation.}
ARC-AGI~\citep{chollet2019measure} and its derivatives (MiniARC,
1D-ARC) operationalise novelty through hand-crafted and synthetic
rule families. We complement those benchmarks by generating
\emph{shapes} rather than \emph{rules}, with a SHA1 hash on rendered
grids ensuring no two grids in the experiment are identical, and by
asking addressing questions rather than transformations.
